\documentclass[11pt]{amsart}
\usepackage{amsmath,amssymb,amsthm}
\usepackage[margin=1in]{geometry}
\usepackage{hyperref}

\newtheorem{theorem}{Theorem}
\newtheorem{lemma}[theorem]{Lemma}
\newtheorem{corollary}[theorem]{Corollary}
\newtheorem{proposition}[theorem]{Proposition}
\theoremstyle{remark}
\newtheorem{remark}[theorem]{Remark}

\newcommand{\T}{\mathbb{T}}
\newcommand{\R}{\mathbb{R}}
\newcommand{\Z}{\mathbb{Z}}
\newcommand{\N}{\mathbb{N}}
\newcommand{\avg}{\frac{1}{2\pi}\int_0^{2\pi}}

\title[A window--Nikolskii inequality]{A window--Nikolskii inequality: sharp one-sided bounds for\\ Littlewood cosine polynomials and dense Chowla sets}
\author{DRAFT --- Chowla cosine project}
\date{October 2026 (draft)}

\begin{abstract}
Let $T(x)=\sum_{j=1}^{n}\operatorname{Re}(b_je^{ijx})$ with $|b_j|\le 1$. We prove a window--Nikolskii inequality relating
$M=-\min_{x\in\R}T(x)$ to the coefficient energy $E=\sum|b_j|^2$ and to the energy in low frequencies. Two consequences
follow. (i) For every choice of phases, $-\min_x\sum_{j=1}^n\cos(jx-\theta_j)\ge 0.28\sqrt n$, and the same lower bound
holds for the maximum; in particular every Littlewood cosine polynomial $\sum_{j=0}^n a_j\cos jx$, $a_j=\pm1$, has
$-\min\ge 0.28\sqrt n-1$. This is the sharp order (Rudin--Shapiro), improving the recently claimed $n^{1/3}/60$.
(ii) Chowla's cosine conjecture holds for dense sets: if $A\subset\N$ is finite with $|A|\ge 0.41\max A$, then
$-\min_x\sum_{a\in A}\cos ax\ge c\sqrt{|A|}$ with an absolute $c>0$; this is the sharp order and improves the
$|A|^{1/4}$ bound recently obtained for $A\subset[1,C|A|]$ in that density range. The same inequality gives
$\sqrt{|A|}$ whenever the frequency profile of $A$ is not front-loaded. We explain why the method cannot be localized
in frequency (it is exactly critical on locally full sets) and show that the excluded front-loaded class contains
sets with $-\min f_A/\sqrt{|A|}\to 1/2$. The 0/1 case of the main inequality and the dense theorem are formalized in
Lean~4 (kernel-checked, standard axioms only).
\end{abstract}

\begin{document}
\maketitle

\section{Introduction}
For a finite set $A$ of positive integers let $f_A(x)=\sum_{a\in A}\cos(ax)$ and $M(A)=-\min_{x\in\R}f_A(x)$.
Chowla conjectured $M(A)\ge c\sqrt{|A|}$. The best general lower bound is polynomial of exponent $1/5$
(Bedert~\cite{Bedert}, Shankar~\cite{Shankar}); the best constructions give $M(A)\le(\frac12+o(1))\sqrt{|A|}$
(Bu et al.~\cite{Bu}). For dense sets, $A\subset[1,C|A|]$, Bu et al.\ prove $M(A)\gg_C|A|^{1/4}$ and remark that this
says nothing about whether $\sqrt{|A|}$ is the true order \cite[Prop.~9.7, Rem.~9.8]{Bu}.

The one-sided Littlewood problem asks the same question for unimodular coefficients: for
$T_n(x)=\sum_{j=1}^n\cos(jx-\theta_j)$, how small can $-\min T_n$ be? Since $\avg T_n=0$ and $\avg T_n^2=n/2$, one has
$\max(\max T_n,-\min T_n)\ge\sqrt{n}/2$ trivially, but each one-sided quantity separately is harder; Erd\'elyi
\cite{Erdelyi} recently claimed $-\min T_n\ge n^{1/3}/60$ and $\max T_n\ge n^{1/3}/60$.

Both questions are instances of one inequality.

\section{The window--Nikolskii inequality}
For $w\in\N$ let $\kappa_k=(1-|k|/(w+1))_+$ and $K_w(y)=\sum_{|k|\le w}\kappa_ke^{iky}\ge0$ (Fej\'er kernel). Put
\[
S_w=\sum_{j=1}^{w}\kappa_j^2=\frac{w(2w+1)}{6(w+1)},\qquad W_w=\sum_{|k|\le w}\kappa_k^2=1+2S_w .
\]
For $L\in\N$ let $C_L$ be any constant such that
\begin{equation}\label{eq:NL}
\avg h^2\le C_L\Big(\avg h\Big)^2\quad\text{for every nonnegative trigonometric polynomial $h$ of degree $\le L$.}
\end{equation}

\begin{lemma}[Nikolskii constants]\label{lem:nik}
\eqref{eq:NL} holds with $C_L=L+1$ and with $C_L=0.811(L+1)+71/300$.
\end{lemma}
\begin{proof}
By Fej\'er--Riesz, $h=|U|^2$ with $U=\sum_{k=0}^{L}u_ke^{ikx}$. Then $h\le(\sum|u_k|)^2\le(L+1)\sum|u_k|^2
=(L+1)\avg h$, and $\avg h^2\le\max h\cdot\avg h$. The second constant is the bound
$\sum_m|(u*\tilde u)(m)|^2\le(0.811(L+1)+71/300)(\sum|u_k|^2)^2$ for the autocorrelation of a sequence of length
$L+1$; its elementary proof (a cumulative-mass estimate) is formalized as \texttt{elementary\_analyticNikolskii}.
\end{proof}

\begin{theorem}\label{thm:main}
Let $b_1,\dots,b_n\in\mathbb C$, $E=\sum_j|b_j|^2$, $T(x)=\sum_{j=1}^n\operatorname{Re}(b_je^{ijx})$, and suppose
$T\ge-M$ on $\R$. Then for every integer $0\le w\le n$, with $L=n+w$,
\[
W_w\Big(M^2+\frac E2\Big)\;\le\;C_L\Big(M^2+\frac12\sum_{j=1}^{w}\kappa_j^2|b_j|^2\Big).
\]
\end{theorem}

\begin{proof}
Let $g=T+M\ge0$. For $x_0\in\R$ the function $h_{x_0}(x)=K_w(x-x_0)g(x)$ is a nonnegative trigonometric polynomial of
degree $\le L$, so $\avg h_{x_0}^2\le C_L I(x_0)^2$ with $I(x_0)=\avg h_{x_0}$. Since $\hat g(0)=M$ and
$\hat g(\pm j)=b_j/2,\bar b_j/2$,
\[
I(x_0)=M+\sum_{j=1}^{w}\kappa_j\operatorname{Re}(b_je^{ijx_0}),\qquad
\avg I(x_0)^2\,dx_0=M^2+\frac12\sum_{j\le w}\kappa_j^2|b_j|^2 .
\]
On the other hand, by Fubini and $\avg K_w(x-x_0)^2dx_0=W_w$ for every $x$,
$\avg\avg h_{x_0}(x)^2\,dx\,dx_0=W_w\avg g^2=W_w(M^2+E/2)$. Averaging the pointwise inequality over $x_0$ gives the claim.
\end{proof}

Only the moduli $|b_j|$ enter: the inequality is blind to phases.

\section{Unimodular and Littlewood cosine polynomials}

\begin{corollary}\label{cor:littlewood}
For every $n\ge1$ and all $\theta_1,\dots,\theta_n\in\R$,
\[
-\min_{x\in\R}\sum_{j=1}^n\cos(jx-\theta_j)\;\ge\;0.28\sqrt n,\qquad
\max_{x\in\R}\sum_{j=1}^n\cos(jx-\theta_j)\;\ge\;0.28\sqrt n .
\]
With the trivial constant $C_L=L+1$ the bound is $\ge0.189\sqrt n$, with asymptotic constant $\sqrt{7/2-2\sqrt3}$.
Consequently every $\sum_{j=0}^n a_j\cos jx$ with $a_j\in\{-1,1\}$ satisfies $-\min\ge0.28\sqrt n-1$ and
$\max\ge0.28\sqrt n-1$.
\end{corollary}

\begin{proof}
Take $|b_j|=1$, so $E=n$ and $\sum_{j\le w}\kappa_j^2|b_j|^2=S_w$. With $C_L=L+1$ and $W_w=1+2S_w$,
Theorem~\ref{thm:main} rearranges to
\[
M^2\;\ge\;\frac{n+S_w(n-w-1)}{2(n+w-2S_w)}\qquad(0\le w\le n-1).
\]
Using $w/3-1/6\le S_w\le w/3$ and $w=\lfloor\theta n\rfloor$ gives $M^2\ge n\,\theta(1-\theta)/(6+2\theta)-O(1)$,
maximized at $\theta=2\sqrt3-3$ with value $(7/2-2\sqrt3)n$. With the second constant of Lemma~\ref{lem:nik} the same
computation (optimum $w\approx0.69n$) gives $M\ge0.2819\sqrt n-o(\sqrt n)$; the exact bound, evaluated for every
$n\le2000$ and estimated beyond, stays above $0.28\sqrt n$. The maximum follows by applying this to $-T$ (phases
$\theta_j+\pi$), and the $\pm1$ statement by adding the constant term $a_0$.
\end{proof}

\begin{remark}
The order is sharp: for $n=2^k$ the Rudin--Shapiro signs give $|\sum_{j<n}\varepsilon_je^{ijx}|\le\sqrt{2n}$, hence
$|\sum_{j=1}^n\varepsilon_{j-1}\cos jx|\le\sqrt{2n}$. Numerically, optimizing the phases gives
$-\min T_n/\sqrt n\approx0.607$ for $8\le n\le384$; chirp phases give $\approx1.17$ and random phases $\approx2.6$.
\end{remark}

\section{Dense Chowla sets}

\begin{corollary}\label{cor:dense}
If $A\subset\N$ is finite and nonempty with $|A|\ge\frac{41}{100}\max A$, then $M(A)\ge\sqrt{|A|}/1000$. More precisely,
if $|A|\ge\delta\max A$ with $\delta>0.4055$, then $M(A)\ge c(\delta)\sqrt{|A|}$ with, for large $|A|$,
$c(0.5)\approx0.064$ and $c(1)\approx0.23$.
\end{corollary}

\begin{proof}[Proof sketch]
Apply Theorem~\ref{thm:main} with $b_a=\mathbf 1_A(a)$ (so $E=|A|$), $X=\max A$, and bound
$\sum_{a\le w}\kappa_a^2\le S_w\approx w/3$. For $w=\theta X$ the inequality becomes
$\frac{2w}{3}\cdot\frac{|A|}2\lesssim 0.811(X+w)\frac{w}{6}+O(M^2X)$, which forces $M^2\gg w$ as soon as
$|A|/X>0.4055(1+\theta)$. The explicit constants are computed in the formal proof.
\end{proof}

\begin{remark}[Front-loading]
More generally, Theorem~\ref{thm:main} gives $M(A)\gg\sqrt{|A|}$ whenever the weighted low-frequency count
$\sum_{a\in A}(1-a/(w+1))_+^2$ is at most $(1-\eta)\frac{2}{3\rho}\frac{w}{X+w}|A|$ for some $w$, where
$\rho=C_L/(L+1)$. Sets violating this at every scale are \emph{front-loaded}, and the class contains very good sets.
Adding a few far elements to a Sidon set $S$ keeps it Sidon, so Bu et al.'s identity
$f_{A(S)}=\rho_S^2-|S|/2-g_S/2$ survives while the profile of $A(S)$ becomes arbitrarily front-loaded. Numerically
(Bu's sets with $|S|=40,80$ plus $2$--$4$ outliers, front-loading factors up to $200$), $M(A)/\sqrt{|A|}$ stays within
$0.03$ of the value for $A(S)$ itself, e.g.\ $0.673$ at $|A|=6972$. In fact, writing $B=\max_x(g_S/2-\rho_S^2)$
($=O(\sqrt{k\log k})$ for Bu et al.'s sets, $k=|S|$) and adding $k_o\le\sqrt{k/2}$ outliers $O$, one has
$g/2-\rho^2\le\min(B,k/2-|\rho_S|^2)+2k_o|\rho_S|+k_o/2\le B+k_o\sqrt{2k}+k_o/2$, so
$M(A(S\cup O))/\sqrt{|A|}\le\frac12+O(\sqrt{\log k/k})+O(k_o/\sqrt k)\to\frac12$ for $k_o=o(\sqrt k)$, while the
front-loading factor is of the order of the outliers' distance ratio. So the front-loading restriction is a limitation of
the method, not of the sets.
\end{remark}

\begin{remark}[Localization is critical]
Let $K\ge0$ be a trigonometric polynomial with $\hat K(0)=1$ and spectrum $\Sigma$, and let
$\rho(\Sigma)=\sup\{\|h\|_2^2/\hat h(0)^2: h\ge0,\ \mathrm{spec}\,h\subset\Sigma\}$. Replacing $h_{x_0}$ by
$h_{x_0}*K$ in the proof of Theorem~\ref{thm:main} yields
\[
M^2\Psi(0)+\tfrac12\sum_{a\in A}\Psi(a)\le\rho(\Sigma)\Big(M^2+\tfrac12\sum_{a\le w}\kappa_a^2\mathbf 1_A(a)\Big),
\qquad \Psi=\hat K^2*\kappa^2 .
\]
If $A$ is full on the positive part of $\Sigma$, the two sides agree to leading order up to the factor
$\|K\|_2^2/\rho(\Sigma)\le1$: no choice of $K$ gives information. The global inequality gains the factor $2/\rho$
because the flat coefficient sequence of $f_A$ is not itself a nonnegative kernel. (This is a leading-order
computation, recorded to explain why the density threshold $0.4055$ cannot be lowered by localizing the window.)
\end{remark}

\section{Formalization}
The following are checked in Lean~4 with Mathlib (no \texttt{sorry}, only the axioms \texttt{propext},
\texttt{Classical.choice}, \texttt{Quot.sound}): the finite Fej\'er--Riesz theorem for every degree; the elementary
Nikolskii constant of Lemma~\ref{lem:nik}; Theorem~\ref{thm:main} for $0/1$ coefficients (\texttt{p31\_elementary});
Corollary~\ref{cor:dense} with constant $1/1000$ (\texttt{dense\_elementary}).
% FILL AFTER AUDIT: complex-coefficient Theorem 2 (p34_main, p34_elementary), Corollary 3 (p34_A: sqrt(n)/4 for every
% n and all phases; p34_B for +-1 polynomials), P33 and the dilation pinch.

\begin{thebibliography}{9}
\bibitem{Bedert} B.~Bedert, \emph{Polynomial bounds for the Chowla cosine problem}, arXiv:2509.05260.
\bibitem{Bu} B.~Bu et al., \emph{A half-constant upper bound for Chowla's cosine problem}, arXiv:2610.17160v1 (2026).
\bibitem{Erdelyi} T.~Erd\'elyi, \emph{Bourgain--Ruzsa type estimates for the maximum and minimum of Littlewood cosine
polynomials on the real line}, arXiv:2607.21877v1 (2026).
\bibitem{Shankar} A.~Shankar, \emph{A log-free $n^{1/5}$ bound for Chowla's cosine problem}, arXiv:2609.05338.
\end{thebibliography}
\end{document}
